% n3_trial_02_actuation_sequence.tex
% Three-panel polymer actuation schematic:
%   Panel 1 — Charge Injection  (+V applied)
%   Panel 2 — Coulomb Repulsion (-V applied)
%   Panel 3 — Relaxation        (Open circuit)
%
% Macro library coverage: none of the published macros (BandDiagram, BellCurve,
% WavyChain, LogLogPlot, IsoBlock, etc.) map to a 2D cantilever mechanism schematic.
% Palette and font style lock are honoured; all drawing is hand-written TikZ.

\documentclass[border=2pt]{standalone}
\usepackage{polymer-paper-preamble}

\begin{document}
\begin{tikzpicture}[
  x=1cm, y=1cm,
  %% shared node font
  every node/.style={font=\sffamily\fontsize{7.5}{9}\selectfont},
  %% filled electron dot (e-)
  edot/.style={
    circle, fill=cGray!90!black,
    inner sep=0pt, minimum size=3pt
  },
  %% electrode block style
  electrode/.style={
    draw=cGray!70!black, fill=cLGray!60,
    line width=0.5pt
  },
  %% force arrow styles
  farrow/.style={
    -{Stealth[length=5pt, width=3.5pt]},
    line width=1.4pt
  },
  %% panel card background
  pcard/.style={
    draw=cLGray!80, fill=cLGray!18,
    rounded corners=4pt, line width=0.5pt
  },
  %% clamp/clip block style
  clamp/.style={
    draw=cGray!80!black, fill=cGray!55,
    line width=0.5pt
  },
  %% ghost dashed outline
  ghost/.style={
    dashed, draw=cGray!55, opacity=0.65, line width=1.0pt
  },
  %% callout annotation box
  callout/.style={
    draw=cTeal!80!black, fill=cTeal!10,
    rounded corners=2pt, inner sep=2pt,
    line width=0.6pt
  },
]

%%==========================================================================
%% Panel dimensions (absolute coordinates, panels laid left-to-right)
%%  Panel 1: x = 0   ... 4.8
%%  Panel 2: x = 5.1 ... 9.9
%%  Panel 3: x = 10.2 ... 15.0
%%  Height: y = 0 ... 7.2
%%==========================================================================

% --- helper: strip width and key dims ---
% strip: thick path, width encoded as line width=10pt
% clamp block top of strip at y = 6.3
% strip anchor x (centre) and electrode x

%%==========================================================================
%% PANEL 1 — Charge Injection
%%==========================================================================
\def\PAx{0}       % panel 1 left x
\def\PAw{4.8}     % panel width
\def\PAmid{2.4}   % panel centre x

% card background
\fill[pcard] (\PAx, 0) rectangle (\PAx+\PAw, 7.2);

% panel title
\node[font=\sffamily\bfseries\fontsize{8.5}{10}\selectfont, anchor=north]
  at (\PAmid, 7.05) {Charge Injection};

% contact clip (GND) label — below clamp block, clear of title
\node[font=\sffamily\fontsize{6.5}{7.8}\selectfont, anchor=north west]
  at (\PAx+0.18, 5.92) {Contact clip (GND)};

% clamp block (fixed support at top of strip)
\fill[clamp] (1.85, 6.00) rectangle (2.95, 6.45);

% strip — slight rightward bend toward +V electrode (briefing §6: bends toward +V)
\draw[line cap=round, line width=10pt, cAmber]
  (2.40, 6.10) .. controls (2.42, 4.50) and (2.58, 3.00) .. (2.75, 1.20);

% strip label
\node[font=\sffamily\fontsize{6.5}{7.8}\selectfont, anchor=east, text=cGray]
  at (1.92, 3.00) {S-rich polymer};

% electrode block (right side, +V)
\fill[electrode] (3.20, 1.20) rectangle (3.85, 5.80);
% electrode label
\node[font=\sffamily\fontsize{7}{8.4}\selectfont, anchor=west]
  at (3.90, 3.50) {+V};

% force arrows (rightward, red — charge injection toward electrode)
\draw[farrow, cRed] (2.55, 4.80) -- (3.15, 4.80);
\draw[farrow, cRed] (2.55, 3.50) -- (3.15, 3.50);
\draw[farrow, cRed] (2.55, 2.20) -- (3.15, 2.20);

% electron dots on strip (positions approximate Bezier t=0.15,0.32,0.50,0.68,0.85)
\node[edot] at (2.42, 5.20) {};
\node[edot] at (2.45, 4.55) {};
\node[edot] at (2.50, 3.85) {};
\node[edot] at (2.57, 3.15) {};
\node[edot] at (2.65, 2.45) {};
% e- text label next to top electron
\node[font=\sffamily\fontsize{6}{7}\selectfont, anchor=west]
  at (2.55, 5.20) {e\textsuperscript{-}};

% bottom caption
\node[font=\sffamily\fontsize{7}{8.4}\selectfont, anchor=north]
  at (\PAmid, -0.10) {+V applied};

%%==========================================================================
%% PANEL 2 — Coulomb Repulsion
%%==========================================================================
\def\PBx{5.1}
\def\PBw{4.8}
\def\PBmid{7.5}   % 5.1 + 2.4

% card background
\fill[pcard] (\PBx, 0) rectangle (\PBx+\PBw, 7.2);

% panel title
\node[font=\sffamily\bfseries\fontsize{8.5}{10}\selectfont, anchor=north]
  at (\PBmid, 7.05) {Coulomb Repulsion};

% clamp block
\fill[clamp] (6.95, 6.00) rectangle (8.05, 6.45);

% strip — bent leftward (Bezier: top fixed at 7.50, bottom deflects to 6.50)
\draw[line cap=round, line width=10pt, cAmber]
  (7.50, 6.10) .. controls (7.45, 4.50) and (7.20, 3.00) .. (6.50, 1.20);

% electrode block (right side, floating / -V)
\fill[electrode] (8.35, 1.20) rectangle (9.00, 5.80);
% electrode label
\node[font=\sffamily\fontsize{7}{8.4}\selectfont, anchor=west]
  at (9.05, 5.40) {floating};
\node[font=\sffamily\fontsize{7}{8.4}\selectfont, anchor=west]
  at (9.05, 3.50) {-V};

% force arrows (leftward, teal — Coulomb repulsion away from electrode)
\draw[farrow, cTeal] (7.30, 4.60) -- (6.65, 4.60);
\draw[farrow, cTeal] (7.10, 3.30) -- (6.45, 3.30);
\draw[farrow, cTeal] (6.90, 2.10) -- (6.25, 2.10);

% electron dots on bent strip (approximate along the Bezier)
\node[edot] at (7.47, 5.20) {};
\node[edot] at (7.42, 4.55) {};
\node[edot] at (7.28, 3.85) {};
\node[edot] at (7.05, 3.15) {};
\node[edot] at (6.80, 2.45) {};
\node[font=\sffamily\fontsize{6}{7}\selectfont, anchor=west]
  at (7.55, 5.20) {e\textsuperscript{-}};

% Coulomb > Maxwell callout annotation — moved below title
\node[callout, font=\sffamily\fontsize{6.5}{7.8}\selectfont]
  at (\PBmid, 6.35) {Coulomb > Maxwell};

% bottom caption
\node[font=\sffamily\fontsize{7}{8.4}\selectfont, anchor=north]
  at (\PBmid, -0.10) {-V applied};

%%==========================================================================
%% PANEL 3 — Relaxation
%%==========================================================================
\def\PCx{10.2}
\def\PCw{4.8}
\def\PCmid{12.6}   % 10.2 + 2.4

% card background
\fill[pcard] (\PCx, 0) rectangle (\PCx+\PCw, 7.2);

% panel title
\node[font=\sffamily\bfseries\fontsize{8.5}{10}\selectfont, anchor=north]
  at (\PCmid, 7.05) {Relaxation};

% clamp block
\fill[clamp] (12.05, 6.00) rectangle (13.15, 6.45);

% ghost — dashed memory of previous (bent) position, same Bezier as panel 2
% shifted into panel 3 coordinate frame (offset +10.2-5.1 = +5.1 in x)
\draw[ghost]
  (12.60, 6.10) .. controls (12.55, 4.50) and (12.30, 3.00) .. (11.60, 1.20);

% "prev." label left of the ghost, clear of the solid strip (C001)
\node[font=\sffamily\fontsize{6}{7}\selectfont, cGray!70, anchor=east]
  at (11.50, 1.35) {prev.};

% strip — partially recovered (less bent than panel 2, slightly right of ghost)
\draw[line cap=round, line width=10pt, cAmber]
  (12.60, 6.10) .. controls (12.58, 4.50) and (12.45, 3.00) .. (11.95, 1.20);

% recovery arrow (curved, from bent position toward upright)
\draw[-{Stealth[length=4.5pt, width=3pt]}, line width=1.2pt, cTeal]
  (11.70, 2.60) .. controls (11.90, 3.20) and (12.10, 3.60) .. (12.30, 3.90);

% electrode block (floating, 0 V)
\fill[electrode] (13.50, 1.20) rectangle (14.15, 5.80);
\node[font=\sffamily\fontsize{7}{8.4}\selectfont, anchor=west]
  at (14.20, 5.40) {floating};

% electron dots on partially-recovered strip
\node[edot] at (12.58, 5.20) {};
\node[edot] at (12.55, 4.55) {};
\node[edot] at (12.47, 3.85) {};
\node[edot] at (12.32, 3.15) {};
\node[edot] at (12.10, 2.45) {};
\node[font=\sffamily\fontsize{6}{7}\selectfont, anchor=west]
  at (12.65, 5.20) {e\textsuperscript{-}};

% 0 V label near electrode
\node[font=\sffamily\fontsize{7}{8.4}\selectfont, anchor=west]
  at (14.20, 3.50) {0 V};

% bottom caption
\node[font=\sffamily\fontsize{7}{8.4}\selectfont, anchor=north]
  at (\PCmid, -0.10) {Open circuit};

\end{tikzpicture}
\end{document}
